\section{Banco de questões — Competências Digitais}

\subsection*{N1 — Fundamento competitivo}
\begin{questao}{\nivel{N1} 02.01 — Cebraspe/Certo ou Errado}O uso de HTTPS protege o canal de comunicação, mas não garante que o conteúdo de um site seja verdadeiro ou que o domínio pertença à entidade que o usuário imagina.\end{questao}
\begin{questao}{\nivel{N1} 02.02 — Cebraspe/Certo ou Errado}O modo de navegação privada impede que o provedor de Internet e a rede corporativa observem os destinos acessados pelo usuário.\end{questao}
\begin{questao}{\nivel{N1} 02.03 — Cebraspe/Certo ou Errado}Compartilhar documento por link ``qualquer pessoa com o link pode editar'' tende a ampliar a superfície de acesso em comparação com compartilhamento nominal para usuários específicos.\end{questao}
\begin{questao}{\nivel{N1} 02.04 — Cebraspe/Certo ou Errado}Em LGPD, consentimento é uma das bases legais possíveis e não deve ser tratado como fundamento obrigatório de todo tratamento de dados pessoais.\end{questao}
\begin{questao}{\nivel{N1} 02.05 — Cebraspe/Certo ou Errado}Dado anonimizado de modo efetivamente irreversível para meios razoáveis disponíveis não é tratado da mesma forma que dado pessoal identificável, ressalvadas hipóteses legais aplicáveis.\end{questao}
\begin{questao}{\nivel{N1} 02.06 — Cebraspe/Certo ou Errado}A LAI determina transparência, mas admite proteção de informações pessoais e outras restrições legítimas previstas em lei.\end{questao}
\begin{questao}{\nivel{N1} 02.07 — Cebraspe/Certo ou Errado}No Marco Civil da Internet, neutralidade de rede deve ser confundida com obrigação de todo tráfego apresentar exatamente a mesma latência e qualidade técnica.\end{questao}
\begin{questao}{\nivel{N1} 02.08 — Cebraspe/Certo ou Errado}Uma senha forte reutilizada em dez serviços continua apresentando risco elevado caso um deles sofra vazamento de credenciais.\end{questao}
\begin{questao}{\nivel{N1} 02.09 — Cebraspe/Certo ou Errado}MFA reduz risco de comprometimento por senha vazada, mas não impede phishing em qualquer modalidade nem substitui análise do contexto de autenticação.\end{questao}
\begin{questao}{\nivel{N1} 02.10 — Cebraspe/Certo ou Errado}Histórico de versões de um documento colaborativo pode apoiar rastreabilidade de alterações, mas não substitui controle de acesso adequado.\end{questao}

\subsection*{N2 — Aplicação}
\begin{questao}{\nivel{N2} 02.11 — FGV/e-mail}Um servidor recebe e-mail urgente supostamente do diretor pedindo transferência de arquivo sigiloso para endereço externo. A melhor conduta inicial é:\begin{enumerate}[label=\Alph*)]\item cumprir imediatamente porque o remetente parece correto;\item confirmar o pedido por canal institucional independente e verificar domínio/destinatário;\item encaminhar a todos os colegas;\item desabilitar o antivírus para abrir o anexo;\item responder com sua senha para confirmar identidade.\end{enumerate}\end{questao}
\begin{questao}{\nivel{N2} 02.12 — FCC/busca}Para localizar um ato normativo específico em domínio oficial, uma estratégia eficiente é:\begin{enumerate}[label=\Alph*)]\item usar expressão exata entre aspas e restringir com \texttt{site:};\item confiar apenas no primeiro resultado patrocinado;\item pesquisar sem termos específicos;\item desabilitar filtros de data;\item usar somente redes sociais.\end{enumerate}\end{questao}
\begin{questao}{\nivel{N2} 02.13 — Cebraspe/Certo ou Errado}Uma fonte secundária pode ser útil para interpretação, mas, quando se busca o texto vigente de uma lei, a fonte oficial/primária deve ser priorizada.\end{questao}
\begin{questao}{\nivel{N2} 02.14 — FGV/permissões}Um aplicativo de lanterna solicita contatos, microfone, localização precisa e acesso contínuo à câmera. A avaliação mais adequada é:\begin{enumerate}[label=\Alph*)]\item conceder tudo, pois apps de loja oficial são sempre seguros;\item questionar proporcionalidade das permissões e negar as não necessárias;\item conceder localização porque qualquer app precisa dela;\item desativar bloqueio de tela;\item compartilhar a conta do dispositivo.\end{enumerate}\end{questao}
\begin{questao}{\nivel{N2} 02.15 — FCC/nuvem}Uma planilha colaborativa deve ser editada por três servidores. A configuração mais segura é:\begin{enumerate}[label=\Alph*)]\item link público com edição;\item compartilhamento nominal com os três usuários e menor permissão necessária;\item publicação na web;\item envio de cópias sem controle de versão;\item conta genérica compartilhada.\end{enumerate}\end{questao}
\begin{questao}{\nivel{N2} 02.16 — Cebraspe/Certo ou Errado}Gravação de videoconferência pode envolver dados pessoais e deve ser tratada conforme finalidade, acesso, retenção e regras institucionais aplicáveis.\end{questao}
\begin{questao}{\nivel{N2} 02.17 — FGV/LGPD}Um órgão trata dados para cumprir obrigação legal. A afirmação correta é:\begin{enumerate}[label=\Alph*)]\item consentimento é obrigatoriamente a única base;\item pode existir base legal diversa do consentimento, conforme a hipótese prevista na LGPD;\item a LGPD não se aplica ao poder público;\item dado pessoal só existe quando sensível;\item todo tratamento exige publicação integral.\end{enumerate}\end{questao}
\begin{questao}{\nivel{N2} 02.18 — FCC/LGPD}Uma empresa decide finalidade e meios essenciais de tratamento; outra processa os dados conforme instruções. Em termos gerais, são, respectivamente:\begin{enumerate}[label=\Alph*)]\item operador e titular;\item controlador e operador;\item encarregado e titular;\item titular e controlador;\item autoridade e encarregado.\end{enumerate}\end{questao}
\begin{questao}{\nivel{N2} 02.19 — Cebraspe/Certo ou Errado}A classificação contratual de uma empresa como ``operadora'' não é suficiente se, na prática, ela toma decisões próprias determinantes sobre finalidade do tratamento.\end{questao}
\begin{questao}{\nivel{N2} 02.20 — FGV/LAI}A divulgação ativa de base pública contém CPF completo de cidadãos sem necessidade demonstrada. A providência adequada é:\begin{enumerate}[label=\Alph*)]\item manter porque transparência sempre supera privacidade;\item avaliar necessidade, regime de acesso e minimização dos dados pessoais;\item remover toda informação pública;\item publicar também endereço residencial;\item ignorar a LGPD.\end{enumerate}\end{questao}
\begin{questao}{\nivel{N2} 02.21 — Cebraspe/Certo ou Errado}Transparência ativa e passiva são formas distintas de implementação do direito de acesso à informação.\end{questao}
\begin{questao}{\nivel{N2} 02.22 — FCC/Marco Civil}A neutralidade de rede está relacionada principalmente a:\begin{enumerate}[label=\Alph*)]\item tratamento isonômico do tráfego nos termos legais, com exceções regulamentadas;\item obrigação de anonimato absoluto;\item proibição de qualquer gestão de rede;\item exclusão de registros de conexão;\item licença de software livre.\end{enumerate}\end{questao}
\begin{questao}{\nivel{N2} 02.23 — FGV/phishing}Um link mostra visualmente \texttt{gov.br}, mas o destino real, ao inspecionar, é domínio semelhante registrado por terceiro. O sinal indica:\begin{enumerate}[label=\Alph*)]\item tentativa potencial de phishing;\item garantia de legitimidade;\item problema de backup;\item neutralidade de rede;\item compressão de URL.\end{enumerate}\end{questao}
\begin{questao}{\nivel{N2} 02.24 — Cebraspe/Certo ou Errado}Encurtadores de URL podem ocultar o destino real e, por isso, merecem cautela em mensagens inesperadas.\end{questao}
\begin{questao}{\nivel{N2} 02.25 — FCC/colaboração}Ao editar simultaneamente documento compartilhado, conflito de versão tende a ser reduzido quando a plataforma oferece:\begin{enumerate}[label=\Alph*)]\item coautoria e histórico de versões;\item cópias independentes por e-mail;\item conta genérica;\item desativação de sincronização;\item link público permanente.\end{enumerate}\end{questao}

\subsection*{N3 — Integração de concurso}
\begin{questao}{\nivel{N3} 02.26 — FGV/assertivas}Considere: I. HTTPS protege confidencialidade/integridade do canal. II. Certificado válido prova que o conteúdo do site é factual. III. Phishing pode ocorrer em site com HTTPS. Assinale:\begin{enumerate}[label=\Alph*)]\item apenas I;\item apenas II;\item I e III;\item II e III;\item I, II e III.\end{enumerate}\end{questao}
\begin{questao}{\nivel{N3} 02.27 — Cebraspe/Certo ou Errado}Uma página de phishing pode usar HTTPS e certificado válido para o domínio fraudulento; portanto, o cadeado não substitui verificação do domínio e do contexto.\end{questao}
\begin{questao}{\nivel{N3} 02.28 — FCC/LGPD}O princípio da necessidade orienta:\begin{enumerate}[label=\Alph*)]\item coleta máxima para usos futuros indefinidos;\item limitação do tratamento ao mínimo necessário para a finalidade;\item eliminação de qualquer dado público;\item uso de consentimento sempre;\item proibição de dados sensíveis.\end{enumerate}\end{questao}
\begin{questao}{\nivel{N3} 02.29 — FGV/LGPD}Um setor coleta telefone, religião e geolocalização para emitir declaração simples que exige apenas nome e matrícula. O principal problema é:\begin{enumerate}[label=\Alph*)]\item excesso de minimização;\item coleta desproporcional à finalidade e necessidade;\item falta de transparência ativa apenas;\item uso de HTTPS;\item ausência de cookies.\end{enumerate}\end{questao}
\begin{questao}{\nivel{N3} 02.30 — Cebraspe/Certo ou Errado}Dados biométricos vinculados a pessoa natural podem integrar categoria de dados pessoais sensíveis, o que exige atenção reforçada ao tratamento.\end{questao}
\begin{questao}{\nivel{N3} 02.31 — FCC/LAI-LGPD}Pedido de acesso envolve documento público com dados pessoais não necessários à compreensão do ato. A solução mais adequada pode incluir:\begin{enumerate}[label=\Alph*)]\item negar sempre o documento inteiro;\item fornecer acesso com ocultação/minimização dos dados protegidos, quando juridicamente cabível;\item publicar todos os dados pessoais;\item exigir consentimento de toda a população em qualquer hipótese;\item ignorar a LAI.\end{enumerate}\end{questao}
\begin{questao}{\nivel{N3} 02.32 — FGV/ECA}A publicação digital de imagem de criança em contexto institucional deve considerar:\begin{enumerate}[label=\Alph*)]\item apenas qualidade da imagem;\item proteção integral, dignidade, privacidade e regras aplicáveis;\item que menores não possuem direitos digitais;\item que conta institucional elimina riscos;\item que toda imagem pública é livre de restrição.\end{enumerate}\end{questao}
\begin{questao}{\nivel{N3} 02.33 — Cebraspe/Certo ou Errado}A proteção de crianças e adolescentes no ambiente digital envolve não apenas dados pessoais, mas também dignidade, integridade e prevenção de exposição danosa.\end{questao}
\begin{questao}{\nivel{N3} 02.34 — FCC/Marco Civil}A guarda e disponibilização de registros de Internet não deve ser tratada como autorização genérica para acesso irrestrito por qualquer agente público, pois existe regime jurídico específico de proteção e requisição.\begin{enumerate}[label=\Alph*)]\item correta;\item incorreta porque registros são sempre públicos;\item incorreta porque nenhum registro pode ser guardado;\item correta apenas para empresas estrangeiras;\item incorreta porque LAI elimina privacidade.\end{enumerate}\end{questao}
\begin{questao}{\nivel{N3} 02.35 — Cebraspe/Certo ou Errado}Em comunicação institucional, mensagem enviada por aplicativo pessoal pode não ser canal adequado para registrar decisão administrativa que precisa de preservação, publicidade interna ou rastreabilidade formal.\end{questao}
\begin{questao}{\nivel{N3} 02.36 — FGV/reunião}Durante videoconferência, um servidor compartilha a tela inteira e exibe notificação com dado sigiloso. O controle preventivo mais adequado inclui:\begin{enumerate}[label=\Alph*)]\item compartilhar apenas janela necessária e desabilitar notificações sensíveis;\item usar resolução maior;\item aumentar volume do microfone;\item gravar automaticamente toda reunião;\item tornar link público.\end{enumerate}\end{questao}
\begin{questao}{\nivel{N3} 02.37 — Cebraspe/Certo ou Errado}Instalar extensão de navegador requer avaliar permissões, porque uma extensão pode ter capacidade de ler ou alterar conteúdo das páginas acessadas.\end{questao}
\begin{questao}{\nivel{N3} 02.38 — FCC/fonte}Uma notícia afirma alteração de lei, mas não cita número nem ato oficial. Para validar, deve-se:\begin{enumerate}[label=\Alph*)]\item compartilhar rapidamente;\item localizar a publicação/ato em fonte oficial e comparar data e conteúdo;\item confiar no número de curtidas;\item usar somente comentários;\item considerar verdadeira se estiver em HTTPS.\end{enumerate}\end{questao}
\begin{questao}{\nivel{N3} 02.39 — FGV/versões}Dois servidores editam cópias locais da mesma planilha e enviam por e-mail ``versão final 3'' e ``versão final nova''. O risco principal é:\begin{enumerate}[label=\Alph*)]\item conflito de versão e perda de fonte única de verdade;\item excesso de colaboração;\item falta de criptografia assimétrica;\item neutralidade;\item ausência de streaming.\end{enumerate}\end{questao}
\begin{questao}{\nivel{N3} 02.40 — Cebraspe/Certo ou Errado}Um gerenciador de senhas pode reduzir reutilização e favorecer credenciais únicas, embora também deva ser protegido por autenticação forte e práticas adequadas.\end{questao}

\subsection*{N4 — Auditoria / avançado}
\begin{questao}{\nivel{N4} 02.41 — Cebraspe/caso integrado}Um setor compartilha por link público uma planilha com nome, CPF e telefone de beneficiários, alegando que ``o link só foi enviado aos servidores''. Julgue: a restrição informal de distribuição do link não equivale a controle de acesso nominal e mantém risco de exposição indevida.\end{questao}
\begin{questao}{\nivel{N4} 02.42 — FGV/incidente}Após envio acidental de dados pessoais a destinatário externo, a primeira resposta adequada inclui:\begin{enumerate}[label=\Alph*)]\item apagar evidências;\item conter o compartilhamento, preservar registros, avaliar escopo/risco e acionar fluxo institucional de incidente;\item ignorar se não houver reclamação;\item publicar os dados para tornar o acesso igualitário;\item desativar todos os e-mails.\end{enumerate}\end{questao}
\begin{questao}{\nivel{N4} 02.43 — FCC/governança}Uma plataforma colaborativa possui 3.000 links públicos antigos sem owner. O controle prioritário é:\begin{enumerate}[label=\Alph*)]\item inventário, revisão de compartilhamentos, owner e política de expiração/minimização;\item aumentar armazenamento;\item criar mais links;\item desabilitar logs;\item migrar para outro fornecedor sem revisar permissões.\end{enumerate}\end{questao}
\begin{questao}{\nivel{N4} 02.44 — Cebraspe/Certo ou Errado}Migrar arquivos para nova plataforma sem revisar permissões pode transportar o mesmo problema de excesso de acesso para o novo ambiente.\end{questao}
\begin{questao}{\nivel{N4} 02.45 — FGV/LAI-LGPD}Um portal publica remuneração de agentes públicos junto com dados bancários e endereço residencial. A melhor análise é:\begin{enumerate}[label=\Alph*)]\item transparência justifica todos os campos;\item deve-se separar informação necessária ao controle público de dados pessoais sem finalidade de divulgação;\item LGPD impede divulgar qualquer remuneração;\item endereço residencial é sempre dado aberto;\item dados bancários devem ser publicados para comprovar pagamento.\end{enumerate}\end{questao}
\begin{questao}{\nivel{N4} 02.46 — Cebraspe/Certo ou Errado}A proteção de dados no setor público deve ser compatibilizada com transparência; tratar uma lei como anuladora automática da outra é simplificação inadequada.\end{questao}
\begin{questao}{\nivel{N4} 02.47 — FCC/phishing}Uma conta institucional com MFA foi comprometida após usuário aprovar notificação inesperada. A conclusão correta é:\begin{enumerate}[label=\Alph*)]\item MFA é inútil;\item fatores e UX de autenticação devem ser analisados, pois MFA também pode sofrer ataques de fadiga/phishing;\item senha não importa;\item HTTPS causou o incidente;\item deve-se remover segundo fator.\end{enumerate}\end{questao}
\begin{questao}{\nivel{N4} 02.48 — FGV/matriz de achado}Auditoria encontra documentos sensíveis compartilhados por links públicos, sem expiração e sem inventário. O achado mais adequado é:\begin{enumerate}[label=\Alph*)]\item ``A nuvem é insegura'';\item ``A ausência de governança de compartilhamento amplia acesso além do necessário e reduz accountability; recomenda-se revisão de links, ownership, expiração e política de permissões'';\item ``Deve-se proibir colaboração'';\item ``Basta mudar a cor da interface'';\item ``Não há risco se os links não aparecem em buscadores''.\end{enumerate}\end{questao}
\begin{questao}{\nivel{N4} 02.49 — Cebraspe/Certo ou Errado}Em avaliação de competência digital, capacidade de verificar fontes e reconhecer limitações de ferramentas é tão relevante quanto saber operar os recursos de software.\end{questao}
\begin{questao}{\nivel{N4} 02.50 — FGV/cidadania digital}Um servidor usa perfil institucional para divulgar opinião partidária como se fosse posição oficial do órgão. O problema principal envolve:\begin{enumerate}[label=\Alph*)]\item apenas largura de banda;\item confusão entre manifestação pessoal e comunicação institucional, com risco à impessoalidade e credibilidade;\item falta de VPN;\item uso de navegador;\item inexistência de arquivo PDF.\end{enumerate}\end{questao}
